%! Sinusoidal Function Context and Data Modeling — Unit 3, Lesson 3.7 — Warm-Up
\documentclass[10pt]{article}
\usepackage{precalculus-article}
\usepackage{precalculus-boxes}
\newcommand{\TallMath}[1]{$\displaystyle #1\rule[-1.4em]{0pt}{3.2em}$}

\begin{document}

\pageheader{Unit 3, Lesson 3.7}{Warm-Up}
\namedateperiod

\vspace{0.12in}

\begin{notesbox}{Spiral Review}

\textbf{1.}\quad The graph below represents one complete period of a sinusoidal
function $g$.

\vspace{6pt}
\begin{center}
\begin{tikzpicture}[scale=0.85]
  \draw[linegray, step=0.5cm] (0,-1.6) grid (6.4, 1.6);
  \draw[->] (0,0) -- (6.6,0) node[right]{\small $x$};
  \draw[->] (0,-1.7) -- (0,1.8) node[above]{\small $g(x)$};
  % Tick marks x
  \foreach \x/\lbl in {1.57/\frac{\pi}{2}, 3.14/\pi, 4.71/\frac{3\pi}{2}, 6.28/2\pi} {
    \draw (\x, 0.06) -- (\x, -0.06);
    \node[below, font=\tiny] at (\x, -0.06) {$\lbl$};
  }
  % Tick marks y
  \foreach \y/\lbl in {-1.5/-3, 1.5/3} {
    \draw (0.06, \y) -- (-0.06, \y);
    \node[left, font=\tiny] at (-0.06, \y) {$\lbl$};
  }
  % Sinusoidal curve: 3sin(x) scaled: amplitude 3 drawn as 1.5 in tikz units
  \draw[navy, thick, domain=0:6.28, samples=100]
    plot (\x, {1.5*sin(\x r)});
  % Midline dashed
  \draw[charcoal, dashed] (0,0) -- (6.6,0);
\end{tikzpicture}
\end{center}

\vspace{6pt}
(a)\quad What is the amplitude of $g$? \writeline

\vspace{4pt}
(b)\quad What is the period of $g$? \writeline

\vspace{4pt}
(c)\quad What is the midline (vertical shift) of $g$? \writeline

\vspace{0.14in}
\textbf{2.}\quad The table below gives values of a periodic function over one period.

\vspace{4pt}
\begin{center}
\renewcommand{\arraystretch}{1.4}
\begin{tabular}{|c|c|c|c|c|c|}
\hline
$t$ & 0 & 3 & 6 & 9 & 12 \\
\hline
$h(t)$ & 2 & 14 & 2 & $-10$ & 2 \\
\hline
\end{tabular}
\end{center}

\vspace{6pt}
(a)\quad What is the maximum value? \underline{\hspace{1.8cm}}
\quad What is the minimum value? \underline{\hspace{1.8cm}}

\vspace{4pt}
(b)\quad Compute the amplitude: $A = \dfrac{\max - \min}{2} = $ \writeline

\vspace{4pt}
(c)\quad Compute the midline: $D = \dfrac{\max + \min}{2} = $ \writeline

\end{notesbox}

\end{document}
